\documentclass[10pt,conference]{IEEEtran}
\usepackage[T1]{fontenc}
\usepackage{lmodern}
\usepackage{amsmath,amssymb}
\usepackage{booktabs}
\usepackage{url}
\usepackage{xspace}
\usepackage[hidelinks]{hyperref}

\input{numbers}

\newcommand{\tool}{PromptABI\xspace}

\begin{document}

\title{Static Detection of Role-Boundary Forgery\\in LLM Chat Templates}

\author{\IEEEauthorblockN{PromptABI contributors}\IEEEauthorblockA{Preprint}}

\maketitle

\begin{abstract}
A chat model receives a conversation as one token sequence. The chat
template, a Jinja program distributed with the model, serializes the messages
and inserts the markers that separate one role's turn from the next. When
message content is interpolated without neutralization and the tokenizer maps a
marker typed by a user to the same token id the template uses, the user can
forge a turn: the model receives exactly the tokens of a genuine assistant or
system turn. We examine this property using the template and
tokenizer files, and we present \tool, a static checker that reports each
attacker-reachable field, each marker it can reproduce, the template line, and
the marker's token ids. \tool combines a bounded symbolic executor with exact
witness paths, a may-taint analysis over the full Jinja AST, and a composition
step that classifies each marker against the real tokenizer. On the
\NModels{} most-downloaded text-generation models on the Hugging Face Hub
(\NTemplates{} distinct templates), an independent oracle that renders with
\texttt{transformers} confirms token-exact forgery of an assistant turn for
every model. \tool flags all of them in a median of \TaintMedianMs{}\,ms per
template; of the markers it classifies as control tokens, \ConfCtrlConfirmed{}
of \ConfCtrlDecided{} are reproduced in user content; 281 have matching genuine
rendered occurrences and 70 have no reference occurrence in the probe conversations. We also evaluate a
serving-side fix: tokenizing message content separately with
\texttt{split\_special\_tokens} removes the forged control tokens in only
\MitBlockedSplit{} of \MitForgeable{} templates, because \MitResidualSplit{}
templates use added tokens that are not marked special, while escaping all
added tokens in content removes them in all \MitBlockedEscape{} templates whose
turn separator contains a control token.
\end{abstract}

\section{Introduction}

Prompt-injection attacks place instructions in data that an LLM-integrated
application passes to a
model~\cite{perez2022ignore,greshake2023indirect,liu2024formalizing}. Most
defenses assume that the application can tell the model which part of its input
is trusted. Structured queries separate instructions from data with reserved
delimiters~\cite{chen2025struq}; preference optimization and instruction
hierarchies train models to respect role
privileges~\cite{chen2024secalign,wallace2024hierarchy}; spotlighting marks
untrusted text~\cite{hines2024spotlighting}. Each relies on role and delimiter
tokens that untrusted text cannot produce.

In open-weight deployments the component that produces those tokens is the
chat template. A Hugging Face model repository ships a Jinja template in
\texttt{tokenizer\_config.json} or \texttt{chat\_template.jinja}, and
\texttt{transformers}~\cite{wolf2020transformers}, vLLM~\cite{kwon2023vllm},
and similar servers render the conversation through it and tokenize the
resulting string. The template writes markers such as
\texttt{<|im\_start|>assistant} around each message's content. If the content
is interpolated verbatim and the tokenizer matches the marker as an added token
anywhere in the string, a user who types the marker produces the template's own
token id, and the model sees a forged turn. This is injection in the sense of
Su and Wassermann: untrusted input changes the parse structure of the generated
program~\cite{su2006essence}. The problem has been raised in issue
trackers~\cite{hf2024issue29279} and used for
jailbreaks~\cite{jiang2025chatbug,zhou2024virtual}, and motivate checking structural exposure across deployed templates.

Our claim is that \emph{role-boundary forgery is a property of the template and
tokenizer files that static analyses can flag per field and marker,
with rendering used to establish token-exact witnesses.} The contributions are:
\begin{itemize}
\item a definition of token-exact role-transition forgery and of per-marker
forgery, separating control-token markers from plain-text markers
(Section~\ref{sec:problem});
\item \tool, a checker that combines bounded symbolic execution,
whole-template taint analysis on the Jinja AST, and tokenizer composition, and
reports field, roles, template line, marker, and token ids
(Sections~\ref{sec:approach} and~\ref{sec:impl});
\item a study of the \NModels{} most-downloaded Hub chat models, together
accounting for \DownloadsMillions{} million downloads in 30 days, with an
independent dynamic oracle, and an evaluation of two serving-side fixes
(Section~\ref{sec:eval}). Scripts, pinned model list, and raw outputs are
released with the tool.
\end{itemize}

\section{Background and Related Work}
\label{sec:related}

\paragraph{Chat templates} A chat template is a Jinja program evaluated in a
sandboxed environment with \texttt{messages}, \texttt{tools}, special-token
variables, and an \texttt{add\_generation\_prompt} flag. The Hugging Face
environment enables \texttt{trim\_blocks} and \texttt{lstrip\_blocks}, adds
\texttt{raise\_exception}, and replaces \texttt{tojson} with
\texttt{json.dumps(\ldots, ensure\_ascii=False)}, which leaves \texttt{<} and
\texttt{>} unescaped. A fast tokenizer matches its added tokens in the rendered
string before the subword model runs. The \texttt{split\_special\_tokens}
option turns this off for tokens marked special, but added tokens can also be
registered as non-special, and those are matched regardless. For plain text,
subword tokenizers such as BPE~\cite{sennrich2016bpe} make a string's ids
depend on its neighbors.

\paragraph{Prompt injection and defenses} Perez and Ribeiro introduced goal
hijacking and prompt leaking~\cite{perez2022ignore}; Greshake et al.\ showed
indirect injection through retrieved content~\cite{greshake2023indirect}; Liu
et al.\ formalized attacks and benchmarked defenses~\cite{liu2024formalizing}.
StruQ separates prompt and data with reserved delimiters and a front end that
keeps the delimiters out of the data~\cite{chen2025struq}; SecAlign and the
instruction hierarchy train role-aware
behavior~\cite{chen2024secalign,wallace2024hierarchy}; spotlighting transforms
untrusted text so the model can recognize it~\cite{hines2024spotlighting}.
\tool checks the precondition these defenses share, that untrusted fields
cannot reproduce role and delimiter tokens, for the templates that open-weight
deployments use.

\paragraph{Special-token and template attacks} ChatBug shows that aligned
models misbehave when inputs depart from the chat format the model was tuned
on~\cite{jiang2025chatbug}. Virtual Context injects special tokens so that a
model treats attacker text as its own output~\cite{zhou2024virtual}. Both
measure model behavior once the structure is forged. \tool reports candidate structural exposures before deployment, with
tokenizer composition and rendering distinguishing their token-level behavior.

\paragraph{Static analysis of templates} Taint analysis is a standard way to
find injection vulnerabilities in web applications, from Pixy for
PHP~\cite{jovanovic2006pixy} to the parse-structure definition of
injection~\cite{su2006essence}. Samuel et al.\ infer context-sensitive
sanitization for web templating languages~\cite{samuel2011autosanitization}.
\tool applies the source, sink, and sanitizer view to Jinja chat templates.
The difference is that the meaning of a sink is decided by a tokenizer: whether
a string is a role boundary depends on the ids the tokenizer gives it.

\paragraph{Structured generation} Grammar-constrained decoding
libraries~\cite{willard2023outlines,dong2024xgrammar,geng2023grammar} constrain
model output to a grammar over tokens. \tool also checks output-side artifacts,
such as stop sequences and grammars, against the tokenizer; this paper
evaluates the input-side role-boundary check.

\section{Problem Definition}
\label{sec:problem}

Let $T$ be a template, $\mathit{render}_T(c)$ the string it produces for
conversation $c$, and $\tau$ a tokenizer mapping strings to id sequences. A
\emph{field} $f$ is a path into the conversation, such as
\texttt{messages[*].content}; it is \emph{attacker-reachable} when an untrusted
party controls its value for some role (we use the user and tool roles). A
\emph{marker} is a string that $T$ emits outside every field: a literal or a
string constant that reaches an output, or a special token the template emits.

\emph{Role-transition forgery.} Let $g$ be the conversation [user: $u$,
assistant: $a$] and $s$ the substring of $\mathit{render}_T(g)$ between $u$
and $a$. $T$ and $\tau$ admit token-exact role-transition forgery when, for the
single-message conversation $h$ = [user: $u\,s\,a$], the ids
$\tau(\mathit{render}_T(h))$ agree with $\tau(\mathit{render}_T(g))$ on every
token before the start of $a$. The model then receives an assistant turn that
the user wrote.

\emph{Marker forgery.} A marker $m$ is forgeable through $f$ when some value of
$f$ places $m$ in the rendered string inside that value's span and $\tau$ gives
this occurrence the same ids as an occurrence the template emits. If $m$ is an
added token with context-independent matching that the tokenizer does not split, $\tau$ maps it to one id in any
context; we call $m$ a \emph{control-token marker}. Otherwise $m$ is a
\emph{plain-text marker}, and equality depends on the surrounding characters.

\section{Approach}
\label{sec:approach}

\paragraph{Marker catalog} \tool collects the strings a template can emit, namely
literal text and string constants that reach an output. Constants used only in
comparisons, filter arguments, or \texttt{raise\_exception} messages are
excluded. From these it extracts markers with a small set of patterns:
angle-bracket and \texttt{<|\ldots|>} tokens, bracketed tags such as
\texttt{[INST]}, code fences, hash headers, plain-text role headers such as
\texttt{Assistant:}, and short compound runs such as
\texttt{<|im\_start|>assistant}. Special tokens declared by the tokenizer that
the template emits are added; declared tokens the template never emits, such as
reserved tokens, are counted separately.

\paragraph{Bounded symbolic execution} For templates in a supported fragment
(simple loops over \texttt{messages}, conditionals on role and index,
\texttt{set}, and a fixed list of filters) \tool enumerates symbolic render
paths up to two messages, one tool, and 128 paths, maps each path's segments to
role regions, and reports a witness for each attacker-controlled expression and
marker: the path, the rendered excerpt with the marker substituted, and the
marker's offset. Outside the fragment or beyond the bounds the executor records
an abstention and keeps the witnesses from the paths it explored.

\paragraph{Whole-template taint analysis} Production templates use
\texttt{namespace()} objects, reversed and filtered loops, macros, string
methods such as \texttt{split} and \texttt{strip}, and
\texttt{raise\_exception}. To cover the full language, \tool parses the template
with Jinja2's own parser, without rendering, and computes a flow-insensitive
may-taint over the AST, iterated to a fixed point. Sources are
\texttt{messages}, \texttt{tools}, and \texttt{documents}; an attribute or item
access on a message record yields a label such as
\texttt{messages[*].content}. Taint flows through assignments (including
namespace attributes and \texttt{set} blocks), loops, slicing, indexing,
concatenation, formatting, conditional expressions, string methods, filters
that return text, and macro calls, whose bodies are analyzed with the call's
argument taints. Direct numeric or boolean outputs are excluded from sinks; their taint is
retained for downstream transformations. Selection filters preserve taint. Replacement preserves candidate taint: deleting a marker can join its
surrounding characters into another occurrence. Jinja's HTML-escaping filters
neutralize directly interpolated string markers that contain \texttt{<},
\texttt{>}, \texttt{\&}, or quotes. \texttt{tojson} is not a sanitizer under the
Hugging Face environment. Each sink records the role guards that dominate it
(\texttt{message.role == 'assistant'} and its negations), so a finding names
the roles whose content reaches the sink, and sinks reachable only from
assistant or system content are reported separately.

\paragraph{Tokenizer composition} For each finding \tool loads
\texttt{tokenizer.json} with the \texttt{tokenizers} library, encodes the
marker, and checks whether it is a single added token with
\texttt{single\_word=False} and no longer overlapping added token, and whether
\texttt{tokenizer\_config.json} sets \texttt{split\_special\_tokens}. The marker
is classified as \emph{control-token}, \emph{split-by-config}, or \emph{text}.
This step reads the tokenizer but never renders the template.

\section{Implementation}
\label{sec:impl}

\tool is a Python package with no required dependencies; the taint analysis
uses \texttt{jinja2} and the composition step uses \texttt{tokenizers}. The
symbolic executor and role-region model are about 2,400 lines, the taint
analysis about 580, and the tokenizer composition about 130.
\texttt{promptabi template-scan DIR} reads a model directory with the file
precedence of \texttt{transformers} (\texttt{chat\_template.jinja}, then
\texttt{chat\_template.json}, then \texttt{tokenizer\_config.json}) and prints
one line per finding, for example
\begin{center}\small\ttfamily
messages[*].content (roles user, line 23) can render\\
'<tool\_call>' [control-token [151657]]
\end{center}
It exits non-zero when user content can render a control-token marker, so it
can gate a CI job. The symbolic check is also the
\texttt{role-boundary-nonforgeability} rule of \texttt{promptabi verify}, which
emits SARIF.

\section{Evaluation}
\label{sec:eval}

\textbf{RQ1} (applicability and cost): on how many deployed templates does each
analysis produce a verdict, and at what cost?
\textbf{RQ2} (model-level detection): does the static verdict agree with
token-exact role-transition forgery observed by rendering?
\textbf{RQ3} (marker-level precision): how many reported markers can user
content reproduce at the token level, and does tokenizer composition predict
which?
\textbf{RQ4} (fixes): does tokenizing message content separately remove the
forgery without changing benign inputs?

\subsection{Setup}

\paragraph{Subjects} The Hub analysis uses checker revision \texttt{4131479};
the synthetic boundary controls evaluate the revised implementation.
On \SnapshotDate{} we listed Hub models with pipeline tag
\texttt{text-generation} by 30-day downloads and kept, in rank order, each
repository that ships a non-empty chat template and a \texttt{tokenizer.json}
at its current commit, until we had \NModels{}; this reached rank
\LastRank{}. Of the \SkipTotal{} repositories skipped, \SkipGGUF{} were
GGUF-only, \SkipNoTemplate{} had no chat template (mostly base models),
\SkipNoTokenizer{} had no \texttt{tokenizer.json}, and \SkipGated{} were gated
without access for our account. The \NModels{} models use \NTemplates{}
distinct templates, because quantized and fine-tuned repositories usually copy
their base model's template; we report per-template counts where duplication
would inflate them. Every file is pinned by commit sha and sha256.

\paragraph{Dynamic oracle} The oracle shares no code with \tool. It loads each
tokenizer with \texttt{transformers} 5.19.0 (falling back to
\texttt{PreTrainedTokenizerFast} when a repository requires remote code),
renders with \texttt{apply\_chat\_template}, and tokenizes with default
settings. For RQ2 it computes the separator $s$ of Section~\ref{sec:problem}
with $u$ = \texttt{Hello there} and a unique string $a$, adding a system
message when a template rejects the first form, and tests token-exact
agreement. For RQ3 it injects each reported marker $m$ into a user message as
\texttt{a~m~b} and compares the ids that cover $m$ with the ids at each place
the template emits $m$ in three genuine conversations: plain, with a system
message, and with a tool call and tool result. If the ids differ it retries
with the two characters that surround a genuine occurrence, as an attacker
would. The outcome is \emph{confirmed}, \emph{confirmed with context},
\emph{confirmed special} (the marker is one of the tokenizer's added tokens
but the probe conversations do not emit it, as with vision or tool-response
tokens), \emph{not confirmed}, or \emph{unobserved} (a plain-text marker the
probe conversations never emit, so there is no reference occurrence).

\subsection{RQ1: Applicability and cost}

The taint analysis parsed all \TaintParsedTpl{} templates with Jinja2's parser
and took a median of \TaintMedianMs{}\,ms per template, \TaintMaxSec{}\,s for
the largest; classifying a marker against a loaded tokenizer took about
\CompMedianMs{}\,ms. The symbolic executor produced witnesses for
\SymFlagModels{} models (\SymFlagTpl{} templates), \SymCompleteTpl{} of them
fully explored within bounds, in a median of \SymMedianSec{}\,s. On the other
\SymAbstainTpl{} templates its fragment does not cover the top-level control
flow, typically a \texttt{namespace()} declaration or a reversed loop over
\texttt{messages}, and it abstains. The taint analysis is therefore the
analysis that scales to the Hub, and the symbolic executor supplies exact
render-path witnesses where its fragment applies. None of the \NModels{}
tokenizer configurations sets \texttt{split\_special\_tokens}.

\subsection{RQ2: Role-transition forgery}

The oracle found token-exact role-transition forgery in all \NModels{} models
and all \NTemplates{} templates. In \SepSpecialTpl{} templates the forged
separator contains at least one added token, for example
\texttt{<|im\_end|>\textbackslash n}\allowbreak\texttt{<|im\_start|>assistant\textbackslash n} in
the ChatML family,
\texttt{<|end|>}\allowbreak\texttt{<|start|>assistant}\allowbreak\texttt{<|channel|>final}\allowbreak\texttt{<|message|>} in gpt-oss, and
\texttt{<\textbar{}Assistant\textbar{}>} (with full-width bars) in DeepSeek-V3;
in the other \MitPlainText{} it is plain text, such as \texttt{[/INST]} in a
Mistral tokenizer where the instruction tags are not added tokens, or
\texttt{Assistant:} in DeepSeek-Coder-V2.

\begin{table}[t]
\centering
\caption{Static verdict against the oracle's token-exact role-transition
forgery. Every subject is forgeable, so the table measures recall.}
\label{tab:rq2}
\begin{tabular}{lrrr}
\toprule
 & flagged & not flagged & recall \\
\midrule
Taint, models & \TaintMTP & \TaintMFN & \TaintMRecall \\
Taint, templates & \TaintTTP & \TaintTFN & \TaintTRecall \\
Symbolic, models & \SymbolicMTP & \SymbolicMFN & \SymbolicMRecall \\
Symbolic, templates & \SymbolicTTP & \SymbolicTFN & \SymbolicTRecall \\
\bottomrule
\end{tabular}
\end{table}

Table~\ref{tab:rq2} shows that the taint analysis flags every model and the
symbolic executor flags the \SymFlagModels{} models whose templates its
fragment covers. Because the sample contains no template that resists the
forgery, model-level precision is not measurable here; RQ3 measures precision
per reported marker instead.

\subsection{RQ3: Marker-level precision}

The oracle checked every marker either analysis reported for
\texttt{messages[*].content}, per distinct template. Of the symbolic executor's
markers, \SrcSymbolicConfirmed{} of \SrcSymbolicDecided{} decided markers were
confirmed (\SrcSymbolicPrecision{}); of the taint analysis's,
\SrcTaintConfirmed{} of \SrcTaintDecided{} (\SrcTaintPrecision{}).
Table~\ref{tab:rq3} splits the taint analysis's markers by tokenizer class.

\begin{table}[t]
\centering
\caption{Oracle outcome for the taint analysis's markers, per distinct
template, by tokenizer-composition class. Precision is confirmed over decided
(all outcomes except unobserved).}
\label{tab:rq3}
\begin{tabular}{lrr}
\toprule
 & control-token & text \\
\midrule
confirmed & \ConfCtrlOutConfirmed & \ConfTextOutConfirmed \\
confirmed with context & \ConfCtrlOutConfirmedContext & \ConfTextOutConfirmedContext \\
confirmed special & \ConfCtrlOutConfirmedSpecial & \ConfTextOutConfirmedSpecial \\
not confirmed & \ConfCtrlOutNotConfirmed & \ConfTextOutNotConfirmed \\
unobserved & \ConfCtrlUnobserved & \ConfTextUnobserved \\
\midrule
precision & \ConfCtrlPrecision & \ConfTextPrecision \\
\bottomrule
\end{tabular}
\end{table}

For control-token markers, 281 findings matched a genuine rendered
occurrence (280 directly and one with context), and 70 additional findings
reproduced an added-token id in user content without a genuine reference
occurrence. These 70 outcomes establish token reachability rather than a
rendered role-boundary comparison. Text markers were confirmed in
\ConfTextConfirmed{} of \ConfTextDecided{} cases, \ConfTextOutConfirmedContext{}
of them only after the attacker reproduces the neighboring characters. The
\ConfTextOutNotConfirmed{} unconfirmed text markers are mostly tags inside tool
instructions, such as \texttt{</IMPORTANT>}, whose boundary BPE pieces merge
with the neighboring characters differently in user content than at the
template's own occurrence, where the next character belongs to a control
token. Tokenizer composition thus gives a severity split: a CI gate on
control-token findings raised no unconfirmed alarm in this sample, and text
findings form a short review list.

\subsection{RQ4: Serving-side fixes}

Both fixes render the conversation with a unique placeholder per message,
tokenize the template's own text normally, and tokenize each message's content
separately. The first tokenizes content with
\texttt{split\_special\_tokens=True}. The second also encodes any added-token
string found in content one character at a time.

For benign conversations both fixes reproduce ordinary tokenization exactly in
\MitBenignIdentical{} of \MitBenignEvaluated{} templates; in the other four,
tokenizing content separately changes a token at a segment boundary. With
\texttt{split\_special\_tokens} alone, every control token of the forged
separator is removed from the user's span in \MitBlockedSplit{} of
\MitForgeable{} templates. In \MitResidualSplit{} templates some separator
tokens survive because they are registered as non-special added tokens (for
example \texttt{<think>} in the Qwen3 family), and in
\MitExactAfterSplitNonPlain{} of them the role
tokens themselves are non-special, so the forged turn remains token-exact;
these are the DeepSeek-V3 and R1 templates and one derivative, whose
\texttt{<\textbar{}User\textbar{}>} and \texttt{<\textbar{}Assistant\textbar{}>}
tokens are added tokens with \texttt{special=false}.
Escaping all added tokens removes every separator control token in all
\MitBlockedEscape{} templates whose separator contains one, leaving no
token-exact forgery among them. The \MitPlainText{} templates with plain-text
separators need a template-side change, because no tokenizer setting
distinguishes a typed \texttt{Assistant:} from the template's.

\subsection{Threats to validity}

The sample is one snapshot of one ranking and omits gated models we could not
access and GGUF-only repositories, whose templates are embedded in the weight
files. Per-model counts are dominated by a few template families, which is why
we also report distinct templates. The oracle renders fixed conversation
shapes: an \emph{unobserved} marker may occur in a conversation we did not
render, and \emph{not confirmed} covers the two contexts we tried. The marker
patterns, in particular the plain-text role-header pattern, were refined while
inspecting this sample, so the recall figures are not a held-out estimate. The taint analysis is heuristic: parsing the full Jinja AST does not establish
soundness for arbitrary transformations. A fixed synthetic stress test
(artifact \texttt{study/audit\_stress.py}) checks eight sanitizer cases:
five single-pass replacement cases reconstruct a removed marker, so replacement
is treated conservatively. Two added-token settings with
\texttt{single\_word=True} also fail context-independent matching and are
classified as text. These diagnostics are synthetic boundary checks. The
property is structural; we measure that the model receives a forged turn, and
prior work measures how models respond to one~\cite{jiang2025chatbug,zhou2024virtual}.

\section{Discussion}

None of the templates in the sample neutralizes message content, so the
protection of role boundaries in these deployments rests entirely on the
serving layer. \texttt{apply\_chat\_template} with default settings passes the
forgery through. Our results suggest two concrete changes. Servers that accept
untrusted message content should tokenize it separately from the template's
text and escape every added token, special or not; RQ4 shows that
\texttt{split\_special\_tokens} alone is insufficient for models that register
role tokens as non-special. Model publishers that use plain-text role headers
should either move to added tokens or escape the header strings in the
template. \tool's \texttt{template-scan} reports both cases and distinguishes
them through the tokenizer class of each finding.

\section{Conclusion}

Template and tokenizer analysis identifies candidate role-boundary
forgery, which rendering can confirm. \tool analyzes all \NTemplates{} distinct templates of
the \NModels{} most-downloaded Hub chat models in milliseconds per template,
flags every model that the rendering oracle shows to be forgeable, and all 281 control-token findings with genuine rendered reference
occurrences match at the token level. The same study shows
that the common serving-side switch, \texttt{split\_special\_tokens}, does not
cover role tokens registered as non-special, and that escaping all added tokens
in content does.

\bibliographystyle{IEEEtran}
\bibliography{references}

\end{document}
